\chapterTocDelegatePageNumber
\chapter{La messa in moto della ruota del Dhamma}

\setTocDelegatedPageNumber
\englishText
\renewcommand{\englishTitle}{La messa in moto della ruota del Dhamma}

\begin{leader}
\soloinstr{Introduzione del solista}

Avendo il Tathāgata raggiunto la suprema e perfetta illuminazione,
 il primo insegnamento che egli espose fu la Ruota del Dhamma insuperabile,
  che una volta messa in moto perfettamente,
   non può essere fermata nel mondo.
    In essa sono dichiarati i due estremi e la Via di Mezzo,
     e la conoscenza e visione purificate nelle Quattro Nobili Verità.
      Questo Sutta, proclamato dal Re del Dhamma,
       che celebra la perfetta illuminazione,
        è conosciuto come "La Messa in moto della Ruota del Dhamma".
         Recitiamo ora questo Sutta come esposizione.

\end{leader}

[Così ho udito.]

In un'occasione il Beato dimorava a Benares, nel parco dei cervi di Isipatana.
Là il Beato si rivolse al gruppo di cinque bhikkhu:
"Bhikkhu, vi sono due estremi che un rinunciante non dovrebbe seguire:
 da una parte la dedizione ai piaceri sensuali, che è bassa, volgare, mondana, ignobile e infruttuosa;
  dall'altra la dedizione all'automortificazione, che è dolorosa, ignobile e infruttuosa.
Evitando questi due estremi, bhikkhu, il Tathāgata ha realizzato la Via di Mezzo che conduce alla visione, alla conoscenza,
 alla calma, alla penetrazione, all'illuminazione e al Nibbāna.
Qual è dunque, bhikkhu, la Via di Mezzo che il Tathāgata ha realizzato,
 che conduce alla visione, alla conoscenza, alla calma, alla penetrazione, all'illuminazione e al Nibbāna?"


\chapterTocSubIndentTrue
\chapter{Dhammacakkappavattana Sutta}

\paliText
\renewcommand{\paliTitle}{Dhammacakkappavattana Sutta}

\begin{leader}
\soloinstr{Introduzione del solista}

\begin{solotwochants}
Anuttaraṁ abhisambodhiṁ & sambujjhitvā tathāgato\\
Pathamaṁ yaṁ adesesi & dhammacakkaṁ anuttaraṁ\\
Sammadeva pavattento & loke appativattiyaṁ\\
Yatthākkhātā ubho antā & paṭipatti ca majjhimā\\
Catūsvāriyasaccesu & visuddhaṁ ñāṇadassanaṁ\\
Desitaṁ dhammarājena & sammāsambodhikittanaṁ\\
Nāmena vissutaṁ suttaṁ & dhammacakkappavattanaṁ\\
Veyyākaraṇapāthena & saṅgītantam bhaṇāma se\\
\end{solotwochants}
\end{leader}

[Evaṁ me sutaṁ]

Ekaṁ samayaṁ bhagavā bārāṇasiyaṁ viharati isipatane migadāye. Tatra kho
bhagavā pañcavaggiye bhikkhū āmantesi:

Dve'me, bhikkhave, antā pabbajitena na sevitabbā: yo cāyaṁ kāmesu
kāma-sukh'allikānuyogo; hīno, gammo, pothujjaniko, anariyo,
anattha-sañhito; yo cāyaṁ atta-kilamathānuyogo; dukkho, anariyo,
anattha-sañhito.

Ete te, bhikkhave, ubho ante anupagamma majjhimā paṭipadā tathāgatena
abhisambuddhā cakkhukaraṇī, ñāṇakaraṇī, upasamāya, abhiññāya,
sambodhāya, nibbānāya saṁvattati.

Katamā ca sā, bhikkhave, majjhimā paṭipadā tathāgatena abhisambuddhā
cakkhukaraṇī ñāṇakaraṇī, upasamāya, abhiññāya, sambodhāya, nibbānāya
saṁvattati.

\clearpage

\englishText
\markboth{\englishTitle}{\rightmark}

Questo è proprio il Nobile Ottuplice Sentiero, vale a dire:
Retta Visione, Retta Intenzione, Retta Parola, Retta Azione, Retto Sostentamento,
Retto Sforzo, Retta Consapevolezza, Retta Concentrazione.

Questa, o Bhikkhu, è la Via di Mezzo realizzata dal Tathāgata,
che genera visione e conoscenza, e conduce alla calma, alla conoscenza superiore,
all'illuminazione, al Nibbāna.

Questa, o bhikkhu, è la Nobile Verità della sofferenza:
La nascita è sofferenza, l'invecchiamento è sofferenza, la morte è sofferenza,
il dolore, il lamento, l'afflizione, la pena e la disperazione sono sofferenza, l'unione con ciò che non piace
è sofferenza, la separazione da ciò che piace è sofferenza, non ottenere ciò che si desidera è sofferenza.
In breve, i cinque aggregati dell'attaccamento sono sofferenza.

Questa, o bhikkhu, è la Nobile Verità dell'origine della sofferenza:

È questa brama che porta a una nuova esistenza, legata al piacere e alla passione,
che si diletta ora qui ora là; cioè, la brama per i piaceri dei sensi,
la brama di esistere e la brama di non esistere.

Questa, o bhikkhu, è la Nobile Verità della cessazione della sofferenza:
È la completa estinzione e cessazione di quella stessa brama, il suo abbandono,
la sua rinuncia, la liberazione e il non attaccamento ad essa.

Questa, o bhikkhu, è la Nobile Verità del sentiero che conduce alla cessazione della sofferenza:


Questo è proprio il Nobile Ottuplice Sentiero, vale a dire: Retta Visione,
Retta Intenzione, Retta Parola, Retta Azione, Retto Sostentamento, Retto Sforzo,
Retta Consapevolezza, Retta Concentrazione.

"Con il pensiero: 'Questa è la Nobile Verità di dukkha',
sorse in me, o bhikkhu, visione, conoscenza, intuizione,
  saggezza, luce, riguardo a cose non conosciute prima."

\clearpage

\paliText
\markboth{\paliTitle}{\rightmark}

Ayam-eva ariyo aṭṭhaṅgiko maggo seyyathīdaṁ:

Sammā-diṭṭhi, sammā-saṅkappo, sammā-vācā, sammā-kammanto, sammā-ājīvo,
sammā-vāyāmo, sammā-sati, sammā-samādhi.

Ayaṁ kho sā, bhikkhave, majjhimā paṭipadā tathāgatena abhisambuddhā
cakkhukaraṇī ñāṇakaraṇī, upasamāya, abhiññāya, sambodhāya, nibbānāya
saṁvattati.

Idaṁ kho pana, bhikkhave, dukkhaṁ ariya-saccaṁ:

Jātipi dukkhā, jarāpi dukkhā, maranampi dukkhaṁ,
soka-parideva-dukkha-domanass'upāyāsāpi dukkhā, appiyehi sampayogo
dukkho, piyehi vippayogo dukkho, yamp'icchaṁ na labhati tampi dukkhaṁ,
saṅkhittena pañcupādānakkhandhā dukkhā.

Idaṁ kho pana, bhikkhave, dukkha-samudayo ariya-saccaṁ:

Yā'yaṁ taṇhā ponobbhavikā nandi-rāga-sahagatā tatra-tatrābhinandinī
seyyathīdaṁ: kāma-taṇhā, bhava-taṇhā, vibhava-taṇhā.

Idaṁ kho pana, bhikkhave, dukkha-nirodho ariya-saccaṁ:

Yo tassā yeva taṇhāya asesa-virāga-nirodho, cāgo, paṭinissaggo, mutti,
anālayo.

Idaṁ kho pana, bhikkhave, dukkha-nirodha-gāminī paṭipadā ariya-saccaṁ:

Ayam-eva ariyo aṭṭhaṅgiko maggo seyyathīdam: sammā-diṭṭhi,
sammā-saṅkappo, sammā-vācā, sammā-kammanto, sammā-ājīvo, sammā-vāyāmo,
sammā-sati, sammā-samādhi.

\enlargethispage{\baselineskip}

[Idaṁ dukkhaṁ] ariya-saccan'ti me bhikkhave, pubbe ananussutesu dhammesu
cakkhuṁ udapādi, ñāṇaṁ udapādi, paññā udapādi, vijjā udapādi, āloko
udapādi.

\clearpage

\englishText
\markboth{\englishTitle}{\rightmark}

‘Con il pensiero: “Questa è la Nobile Verità di dukkha, e questo dukkha
deve essere compreso”, sorse in me, o bhikkhu, visione, conoscenza,
intuizione, saggezza, luce, riguardo a cose non conosciute prima.

‘Con il pensiero: “Questa è la Nobile Verità di dukkha, e questo dukkha
è stato compreso”, sorse in me, o bhikkhu, visione, conoscenza,
intuizione, saggezza, luce, riguardo a cose non conosciute prima.

‘Con il pensiero: “Questa è la Nobile Verità della causa del dukkha”,
sorse in me, o bhikkhu, visione, conoscenza, intuizione, saggezza, luce,
riguardo a cose non conosciute prima.

‘Con il pensiero: “Questa è la Nobile Verità della causa del dukkha, e
questa causa del dukkha deve essere abbandonata”, sorse in me, o bhikkhu,
visione, conoscenza, intuizione, saggezza, luce, riguardo a cose non conosciute
prima.
before.

‘Con il pensiero: “Questa è la Nobile Verità della causa del dukkha, e
questa causa del dukkha è stata abbandonata”, sorse in me, o bhikkhu,
visione, conoscenza, intuizione, saggezza, luce, riguardo a cose non conosciute
prima.

‘Con il pensiero: “Questa è la Nobile Verità della cessazione del dukkha”,
sorse in me, o bhikkhu, visione, conoscenza, intuizione, saggezza, luce,
riguardo a cose non conosciute prima.

‘Con il pensiero: “Questa è la Nobile Verità della cessazione del dukkha,
e questa cessazione del dukkha deve essere realizzata”, sorse in me,
o bhikkhu, visione, conoscenza, intuizione, saggezza, luce, riguardo a cose
non conosciute prima.

‘Con il pensiero: “Questa è la Nobile Verità della cessazione del dukkha,
e questa cessazione del dukkha è stata realizzata”, sorse in me,
o bhikkhu, visione, conoscenza, intuizione, saggezza, luce, riguardo a cose
non conosciute prima.

\clearpage

\paliText
\markboth{\paliTitle}{\rightmark}

Taṁ kho pan'idaṁ dukkhaṁ ariya-saccaṁ pariññeyyan'ti me bhikkhave, pubbe
ananussutesu dhammesu cakkhuṁ udapādi, ñāṇaṁ udapādi, paññā udapādi,
vijjā udapādi, āloko udapādi.

Taṁ kho pan'idaṁ dukkhaṁ ariya-saccaṁ pariññātan'ti me bhikkhave, pubbe
ananussutesu dhammesu cakkhuṁ udapādi, ñāṇaṁ udapādi, paññā udapādi,
vijjā udapādi, āloko udapādi.

Idaṁ dukkha-samudayo ariya-saccan'ti me bhikkhave, pubbe ananussutesu
dhammesu cakkhuṁ udapādi, ñāṇaṁ udapādi, paññā udapādi, vijjā udapādi,
āloko udapādi.

Taṁ kho pan'idaṁ dukkha-samudayo ariyasaccaṁ pahātabban'ti me bhikkhave,
pubbe ananussutesu dhammesu cakkhuṁ udapādi, ñāṇaṁ udapādi, paññā
udapādi, vijjā udapādi, āloko udapādi.

Taṁ kho pan'idaṁ dukkha-samudayo ariya-saccaṁ pahīnan'ti me bhikkhave, pubbe
ananussutesu dhammesu cakkhuṁ udapādi, ñāṇaṁ udapādi, paññā udapādi,
vijjā udapādi, āloko udapādi.

Idaṁ dukkha-nirodho ariya-saccan'ti me bhikkhave, pubbe ananussutesu
dhammesu cakkhuṁ udapādi, ñāṇaṁ udapādi, paññā udapādi, vijjā udapādi,
āloko udapādi.

Taṁ kho pan'idaṁ dukkha-nirodho ariya-saccaṁ sacchikātabban'ti me bhikkhave,
pubbe ananussutesu dhammesu cakkhuṁ udapādi, ñāṇaṁ udapādi, paññā
udapādi, vijjā, udapādi āloko udapādi.

Taṁ kho pan'idaṁ dukkha-nirodho ariya-saccaṁ sacchikatan'ti me bhikkhave,
pubbe ananussutesu dhammesu cakkhuṁ udapādi, ñāṇaṁ udapādi, paññā
udapādi, vijjā udapādi, āloko udapādi.

\clearpage

\englishText
\markboth{\englishTitle}{\rightmark}

‘Con il pensiero: “Questa è la Nobile Verità del sentiero che conduce alla
cessazione del dukkha”, sorse in me, o bhikkhu, visione, conoscenza,
intuizione, saggezza, luce, riguardo a cose non conosciute prima.

‘Con il pensiero: “Questa Nobile Verità del sentiero che conduce alla cessazione
del dukkha deve essere sviluppata”, sorse in me, o bhikkhu, visione,
conoscenza, intuizione, saggezza, luce, riguardo a cose non conosciute prima.

‘Con il pensiero: “Questa Nobile Verità del sentiero che conduce alla cessazione
del dukkha è stata sviluppata”, sorse in me, o bhikkhu, visione,
conoscenza, intuizione, saggezza, luce, riguardo a cose non conosciute prima.

‘Finché, o bhikkhu, la mia conoscenza e visione della realtà riguardo
a queste Quattro Nobili Verità, nelle loro tre fasi e dodici aspetti, non
fu pienamente chiara per me, non dichiarai al mondo degli spiriti,
dei demoni e degli dei, con i suoi cercatori e saggi, esseri celesti e umani,
la realizzazione dell'incomparabile, perfetta illuminazione.

‘Ma quando, o bhikkhu, la mia conoscenza e visione della realtà riguardo a queste
Quattro Nobili Verità, nelle loro tre fasi e dodici aspetti, fu pienamente
chiara per me, dichiarai al mondo degli spiriti, dei demoni e degli dei, con
i suoi cercatori e saggi, esseri celesti e umani, di aver realizzato
l'incomparabile, perfetta illuminazione.


‘Sorse la conoscenza e la visione: “Incrollabile è la mia liberazione; questa è
l'ultima nascita, non ci sarà più un rinnovamento dell'essere.””\thinspace ’

Così parlò il Beato. Lieto nel cuore, il gruppo di cinque bhikkhu
approvò le parole del Beato.

\clearpage

\paliText
\markboth{\paliTitle}{\rightmark}

Idaṁ dukkha-nirodha-gāminī paṭipadā ariya-saccan'ti me bhikkhave, pubbe
ananussutesu dhammesu cakkhuṁ udapādi, ñāṇaṁ udapādi, paññā udapādi,
vijjā udapādi, āloko udapādi.

Taṁ kho pan'idaṁ dukkha-nirodha-gāminī paṭipadā ariya-saccaṁ bhāvetabban'ti
me bhikkhave, pubbe ananussutesu dhammesu cakkhuṁ udapādi, ñāṇaṁ
udapādi, paññā udapādi, vijjā udapādi, āloko udapādi.

Taṁ kho pan'idaṁ dukkha-nirodha-gāminī paṭipadā ariya-saccaṁ bhāvitan'ti me
bhikkhave, pubbe ananussutesu dhammesu cakkhuṁ udapādi, ñāṇaṁ udapādi,
paññā udapādi, vijjā udapādi, āloko udapādi.

[Yāva kīvañca me bhikkhave,] imesu catūsu ariya-saccesu evan-ti-parivaṭṭaṁ
dvādas'ākāraṁ yathā-bhūtaṁ ñāṇa-dassanaṁ na suvisuddhaṁ ahosi, n'eva tāv'āhaṁ
bhikkhave, sadevake loke samārake sabrahmake sassamaṇa-brāhmaṇiyā pajāya
sadeva-manussāya anuttaraṁ sammā-sambodhiṁ abhisambuddho paccaññāsiṁ.

Yato ca kho me bhikkhave, imesu catūsu ariya-saccesu evan-ti-parivaṭṭaṁ
dvādas'ākāraṁ yathā-bhūtaṁ ñāṇa-dassanaṁ suvisuddham ahosi, ath'āham
bhikkhave, sadevake loke samārake sabrahmake sassamaṇa-brāhmaṇiyā pajāya
sadeva-manussāya anuttaraṁ sammā-sambodhiṁ abhisambuddho paccaññāsiṁ.

Ñāṇañca pana me dassanaṁ udapādi, akuppā me vimutti ayam-antimā jāti,
natthi dāni punabbhavo'ti.

Idam-avoca bhagavā. Attamanā pañcavaggiyā bhikkhū bhagavato bhāsitaṁ
abhinanduṁ.

\clearpage

\englishText
\markboth{\englishTitle}{\rightmark}

Mentre questa esposizione procedeva, la visione immacolata e senza macchia del
Dhamma apparve al Venerabile Koṇḍañña ed egli seppe: ‘Tutto ciò
che ha la natura di sorgere ha la natura di cessare.’

Quando il Beato mise in moto la Ruota del Dhamma, i
deva legati alla Terra proclamarono con una sola voce,

‘L'incomparabile Ruota del Dhamma è stata messa in moto dal Beato
nel santuario dei cervi a Isipatana, vicino a Benares, e nessun cercatore,
bramino, essere celeste, demone, dio o qualsiasi altro essere al mondo
può fermarla.’

Avendo udito ciò che i deva legati alla Terra dissero, i deva dei Quattro Grandi
Re proclamarono con una sola voce\ldots

Avendo udito ciò che i deva dei Quattro Grandi Re dissero, i deva dei
Trentatré proclamarono con una sola voce\ldots

Avendo udito ciò che i deva dei Trentatré dissero, gli Yāma deva
proclamarono con una sola voce\ldots

Avendo udito ciò che gli Yāma deva dissero, i Deva della Gioia proclamarono
con una sola voce\ldots

Avendo udito ciò che i Deva della Gioia dissero, i Deva che si dilettano
nel creare, proclamarono con una sola voce\ldots

Avendo udito ciò che i Deva che si dilettano nel creare dissero, i Deva che
si dilettano nelle creazioni altrui proclamarono con una sola voce\ldots

Avendo udito ciò che i Deva che si dilettano nelle creazioni altrui dissero,
gli dei Brahma proclamarono con una sola voce,

‘L'incomparabile Ruota del Dhamma è stata messa in moto dal Beato
nel santuario dei cervi a Isipatana, vicino a Benares, e nessun cercatore,
bramino, essere celeste, demone, dio o qualsiasi altro essere al mondo
può fermarla.’

\clearpage

\paliText
\markboth{\paliTitle}{\rightmark}

Imasmiñca pana veyyākaraṇasmiṁ bhaññamāne āyasmato koṇḍaññassa virajaṁ
vītamalaṁ dhammacakkhuṁ udapādi: yaṁ kiñci samudaya-dhammaṁ sabban-taṁ
nirodha-dhamman'ti.

[Pavattite ca bhagavatā] dhammacakke bhummā devā saddamanussāvesuṁ:

Etaṁ bhagavatā bārāṇasiyaṁ isipatane migadāye anuttaraṁ dhammacakkaṁ
pavattitaṁ appaṭivattiyaṁ samaṇena vā brāhmaṇena vā devena vā mārena vā
brahmunā vā kenaci vā lokasmin'ti.

Bhummānaṁ devānaṁ saddaṁ sutvā, cātummahārājikā devā
saddamanussāvesuṁ\ldots

Cātummahārājikānaṁ devānaṁ saddaṁ sutvā, tāvatiṁsā devā
saddamanussāvesuṁ\ldots

Tāvatiṁsānaṁ devānaṁ saddaṁ sutvā, yāmā devā saddamanussāvesuṁ\ldots

Yāmānaṁ devānaṁ saddaṁ sutvā, tusitā devā saddamanussāvesuṁ\ldots

Tusitānaṁ devānaṁ saddaṁ sutvā, nimmānaratī devā saddamanussāvesum\ldots

Nimmānaratīnaṁ devānaṁ saddaṁ sutvā, paranimmitavasavattī devā
saddamanussāvesuṁ\ldots

Paranimmitavasavattīnaṁ devānaṁ saddaṁ sutvā, brahmakāyikā devā
saddamanussāvesuṁ:

Etaṁ bhagavatā bārāṇasiyaṁ isipatane migadāye anuttaraṁ dhammacakkaṁ
pavattitaṁ appaṭivattiyaṁ samaṇena vā brāhmaṇena vā devena vā mārena vā
brahmunā vā kenaci vā lokasmin'ti.

\clearpage

\englishText
\markboth{\englishTitle}{\rightmark}

Così in un momento, un istante, un lampo, la notizia della Messa in Moto della
Ruota del Dhamma si diffuse fino al mondo di Brahma, e il
sistema universale decemila volte tremò, si scosse e vibrò, e uno
splendore sublime e sconfinato, che superava il potere dei deva, apparve sulla
terra.

Allora il Beato pronunciò l'enunciazione,

‘In verità, Koṇḍañña ha compreso, Koṇḍañña ha compreso!’ Fu così
che il Venerabile Koṇḍañña ricevette il nome di Aññā-Koṇḍañña: ‘Koṇḍañña Che
Comprende.’

Così termina il discorso sulla Messa in Moto della Ruota del Dhamma.

\clearpage

\paliText
\markboth{\paliTitle}{\rightmark}

Iti'ha tena khaṇena, tena muhuttena, yāva brahmalokā saddo abbhuggacchi.
Ayañca dasa-sahassī lokadhātu saṅkampi sampakampi sampavedhi, appamāṇo ca
oḷāro obhāso loke pāturahosi atikkammeva devānaṁ devānubhāvaṁ.

Atha kho bhagavā udānaṁ udānesi:

Aññāsi vata bho koṇḍañño, aññāsi vata bho koṇḍañño'ti. Iti hidaṁ āyasmato
koṇḍaññassa aññā-koṇḍañño tveva nāmaṁ ahosī'ti.

Dhammacakkappavattana-suttaṁ niṭṭhitaṁ.

\chapterTocDelegatePageNumber
\chapter{La caratteristica del non-sé}

\setTocDelegatedPageNumber
\englishText
\renewcommand{\englishTitle}{Anatta-lakkhaṇa Sutta}

\begin{leader}
\soloinstr{Introduzione del solista}


Tutti gli esseri dovrebbero sforzarsi di comprendere la caratteristica del
non-sé, che offre una liberazione ineguagliabile dalla visione e percezione del sé,
come insegnato dal Buddha supremo.
Questo insegnamento è dato affinché coloro che meditano sulle realtà esperibili
possano giungere a una perfetta comprensione;
per lo sviluppo di una perfetta comprensione di questi fenomeni
e per l'indagine di tutti i momenti mentali inquinati.
La conseguenza di questa pratica è la liberazione totale; quindi, desiderosi di
portare avanti questo insegnamento con il suo grande beneficio,
recitiamo ora questo Sutta.


\end{leader}

[Così ho udito.]

In un'occasione il Beato dimorava a Benares, nel parco dei cervi di Isipatana.
Lì si rivolse al gruppo di cinque bhikkhu:

‘La forma, o bhikkhu, è non-sé. Se, o bhikkhu, la forma fosse il sé, allora la forma
non porterebbe all'afflizione, e si potrebbe dire riguardo alla
forma: “Che la mia forma sia così, che la mia forma non sia così”. Ma poiché,
o bhikkhu, la forma è non-sé, la forma porta quindi all'afflizione, e non
si può dire riguardo alla forma: “Che la mia forma sia così, che la mia forma
non sia così”.

‘La sensazione è non-sé. Se, o bhikkhu, la sensazione fosse il sé, la sensazione
non porterebbe all'afflizione, e si potrebbe dire riguardo alla
sensazione: “Che la mia sensazione sia così, che la mia sensazione non sia così”. Ma
poiché, o bhikkhu, la sensazione è non-sé, la sensazione porta quindi
all'afflizione, e non si può dire riguardo alla sensazione: “Che la mia
sensazione sia così, che la mia sensazione non sia così”.


\chapterTocSubIndentTrue
\chapter{Anatta-lakkhaṇa Sutta}

\paliText
\renewcommand{\paliTitle}{Anatta-lakkhaṇa Sutta}

\begin{leader}
\soloinstr{Introduzione del solista}

{\setlength{\tabcolsep}{0.9em}
\begin{solotwochants}
Yantaṁ sattehi dukkhena & ñeyyaṁ anattalakkhaṇaṁ\\
Attavādattasaññāṇaṁ  & sammadeva vimocanaṁ\\
Sambuddho taṁ pakāsesi & diṭṭhasaccāna yoginaṁ\\
Uttariṁ paṭivedhāya & bhāvetuṁ ñāṇamuttamaṁ\\
Yantesaṁ diṭṭhadhammānam & ñāṇenupaparikkhataṁ\\
Sabbāsavehi cittāni & vimucciṁsu asesato\\
Tathā ñāṇānussārena & sāsanaṁ kātumicchataṁ\\
Sādhūnaṁ atthasiddhatthaṁ & taṁ suttantaṁ bhaṇāma se\\
\end{solotwochants}
}
\end{leader}

[Evaṁ me sutaṁ]

Ekaṁ samayaṁ bhagavā bārāṇasiyaṁ viharati isipatane migadāye. Tatra kho
bhagavā pañcavaggiye bhikkhū āmantesi:

Rūpaṁ bhikkhave anattā, rūpañca hidaṁ bhikkhave attā abhavissa, nayidaṁ rūpaṁ
ābādhāya saṁvatteyya, labbhetha ca rūpe, evaṁ me rūpaṁ hotu, evaṁ me rūpaṁ mā
ahosī'ti. Yasmā ca kho bhikkhave rūpaṁ anattā, tasmā rūpaṁ ābādhāya saṁvattati,
na ca labbhati rūpe, evaṁ me rūpaṁ hotu, evaṁ me rūpaṁ mā ahosī'ti.

Vedanā anattā, vedanā ca hidaṁ bhikkhave attā abhavissa, nayidaṁ vedanā ābādhāya
saṁvatteyya, labbhetha ca vedanāya, evaṁ me vedanā hotu, evaṁ me vedanā mā
ahosī'ti. Yasmā ca kho bhikkhave vedanā anattā, tasmā vedanā ābādhāya
saṁvattati, na ca labbhati vedanāya, evaṁ me vedanā hotu, evaṁ me vedanā mā
ahosī'ti.

\clearpage

\englishText
\markboth{\englishTitle}{\rightmark}

‘La percezione è non-sé. Se, o bhikkhu, la percezione fosse il sé, la percezione
non porterebbe all'afflizione, e si potrebbe dire riguardo alla percezione: “Che la mia
percezione sia così, che la mia percezione non sia così”. Ma poiché, o bhikkhu, la
percezione è non-sé, la percezione porta quindi all'afflizione, e non si può dire
riguardo alla percezione: “Che la mia percezione sia così, che la mia percezione non sia
così”.

‘Le formazioni mentali sono non-sé. Se, o bhikkhu, le formazioni mentali fossero il
sé, le formazioni mentali non porterebbero all'afflizione, e si potrebbe dire riguardo
alle formazioni mentali: “Che le mie formazioni mentali siano così, che le mie
formazioni mentali non siano così”. Ma poiché, o bhikkhu, le formazioni mentali non
sono il sé, le formazioni mentali portano quindi all'afflizione, e non si può dire
riguardo alle formazioni mentali: “Che le mie formazioni mentali siano così, che le mie
formazioni mentali non siano così”.

‘La coscienza è non-sé. Se, o bhikkhu, la coscienza fosse il sé, la coscienza non
porterebbe all'afflizione, e si potrebbe dire riguardo alla coscienza: “Che la mia
coscienza sia così, che la mia coscienza non sia così”. Ma poiché, o bhikkhu, la
coscienza è non-sé, la coscienza porta quindi all'afflizione, e non si può dire
riguardo alla coscienza: “Che la mia coscienza sia così, che la mia coscienza non sia
così”.


Cosa pensate di questo, bhikkhu? La forma è permanente o
impermanente?’


‘Impermanente, Venerabile Signore ’.

‘Ma ciò che è impermanente è doloroso o piacevole?’

‘Doloroso, Venerabile Signore ’.

‘Ma è giusto considerare ciò che è impermanente, doloroso, dinatura mutevole, come
“Questo è mio, io sono questo, questo è il mio sé”?’

‘Non è giusto, Venerabile Signore .’

\clearpage

\paliText
\markboth{\paliTitle}{\rightmark}

Saññā anattā, saññā ca hidaṁ bhikkhave attā abhavissa, nayidaṁ saññā ābādhāya
saṁvatteyya, labbhetha ca saññāya, evaṁ me saññā hotu, evaṁ me saññā mā
ahosī'ti. Yasmā ca kho bhikkhave saññā anattā, tasmā, saññā ābādhāya saṁvattati,
na ca labbhati saññāya, evaṁ me saññā hotu, evaṁ me saññā mā ahosī'ti.

Saṅkhārā anattā, saṅkhārā ca hidaṁ bhikkhave attā abhavissaṁsu, nayidaṁ saṅkhārā
ābādhāya saṁvatteyyuṁ, labbhetha ca saṅkhāresu, evaṁ me saṅkhārā hontu, evaṁ me
saṅkhārā mā ahesun'ti. Yasmā ca kho bhikkhave saṅkhārā anattā, tasmā saṅkhārā
ābādhāya saṁvattanti, na ca labbhati saṅkhāresu, evaṁ me saṅkhārā hontu, evaṁ me
saṅkhārā mā ahesun'ti.

Viññāṇaṁ anattā, viññāṇañca hidaṁ bhikkhave attā abhavissa, nayidaṁ viññānam
ābādhāya saṁvatteyya, labbhetha ca viññāne evaṁ me viññāṇaṁ hotu, evaṁ me
viññāṇaṁ mā ahosī'ti. Yasmā ca kho bhikkhave viññāṇaṁ anattā, tasmā viññāṇaṁ
ābādhāya saṁvattati, na ca labbhati viññāne, evaṁ me viññāṇaṁ hotu, evaṁ me
viññāṇaṁ mā ahosī'ti.

[Taṁ kiṁ maññatha bhikkhave,] rūpam niccaṁ vā aniccaṁ vā'ti.

Aniccaṁ bhante.

Yam panāniccaṁ, dukkhaṁ vā taṁ sukhaṁ vā'ti.

Dukkhaṁ bhante.

Yam panāniccaṁ dukkhaṁ viparināma-dhammaṁ, kallaṁ nu taṁ samanupassituṁ,
etaṁ mama, esoham'asmi, eso me attā'ti.

No hetaṁ bhante.

\clearpage

\englishText
\markboth{\englishTitle}{\rightmark}

Cosa pensate di questo, bhikkhus? La sensazione è permanente o
impermanente?’


‘Impermanente, Venerabile Signore ’.

‘Ma ciò che è impermanente è doloroso o piacevole?’ 

‘Doloroso, Venerabile Signore ’.

‘Ma è giusto considerare ciò che è impermanente, doloroso, dinatura mutevole, come
“Questo è mio, io sono questo, questo è il mio sé”?’

‘Non è giusto, Venerabile Signore .’

Cosa pensate di questo, bhikkhus? La percezione è permanente o
impermanente?’

‘Impermanente, Venerabile Signore ’.

‘Ma ciò che è impermanente è doloroso o piacevole?’

‘Doloroso, Venerabile Signore ’.


‘Ma è giusto considerare ciò che è impermanente, doloroso, dinatura mutevole, come
“Questo è mio, io sono questo, questo è il mio sé”?’

‘Non è giusto, Venerabile Signore .’

Cosa pensate di questo, bhikkhus? Le formazioni mentali sono permanenti o
impermanenti?’

‘Impermanenti, Venerabile Signore ’.

‘Ma ciò che è impermanente è doloroso o piacevole?’ 
‘Dolorose, Venerabile Signore ’.
‘Ma è giusto considerare ciò che è impermanente, doloroso, dinatura mutevole, come
“Questo è mio, io sono questo, questo è il mio sé”?’

‘Non è giusto, Venerabile Signore .’

\clearpage

\paliText
\markboth{\paliTitle}{\rightmark}

Taṁ kiṁ maññatha bhikkhave, vedanā niccā vā aniccā vā'ti.

Aniccā bhante.

Yam panāniccaṁ, dukkhaṁ vā taṁ sukhaṁ vā'ti.

Dukkhaṁ bhante.

Yam panāniccaṁ dukkhaṁ viparināma-dhammaṁ, kallaṁ nu taṁ samanupassituṁ,
etaṁ mama, esoham'asmi, eso me attā'ti.

No hetaṁ bhante.

Taṁ kiṁ maññatha bhikkhave, saññā niccā vā aniccā vā'ti.

Aniccā bhante.

Yam panāniccaṁ, dukkhaṁ vā taṁ sukhaṁ vā'ti.

Dukkhaṁ bhante.

Yam panāniccaṁ dukkhaṁ viparināma-dhammaṁ, kallaṁ nu taṁ samanupassituṁ,
etaṁ mama, esoham'asmi, eso me attā'ti.

No hetaṁ bhante.

Taṁ kiṁ maññatha bhikkhave, saṅkhārā niccā vā aniccā vā'ti.

Aniccā bhante.

Yam panāniccaṁ, dukkhaṁ vā taṁ sukhaṁ vā'ti.

Dukkhaṁ bhante.

Yam panāniccaṁ dukkhaṁ viparināma-dhammaṁ, kallaṁ nu taṁ samanupassituṁ,
etaṁ mama, esoham'asmi, eso me attā'ti.

No hetaṁ bhante.

\clearpage

\englishText
\markboth{\englishTitle}{\rightmark}

Cosa pensate di questo, bhikkhus? La coscienza è permanente o
impermanente?’


‘Impermanente, Venerabile Signore ’.

‘Ma ciò che è impermanente è doloroso o piacevole?’ 

‘Doloroso, Venerabile Signore ’.

‘Ma è giusto considerare ciò che è impermanente, doloroso, dinatura mutevole, come
“Questo è mio, io sono questo, questo è il mio sé”?’

‘Non è giusto, Venerabile Signore .’

‘Perciò, o bhikkhu, qualsiasi forma ci sia, passata, futura, presente,
interna o esterna, grossolana o sottile, inferiore o superiore, che sia
lontana o vicina, ogni forma dovrebbe, per mezzo della retta saggezza, essere vista così com'è
realmente: “Questa non è mia, io non sono questo, questo non è il mio sé”.

‘Qualsiasi sensazione ci sia, passata, futura, presente, interna o
esterna, grossolana o sottile, inferiore o superiore, che sia lontana o
vicina, ogni sensazione dovrebbe, per mezzo della retta saggezza, essere vista così com'è realmente:
“Questa non è mia, io non sono questo, questo non è il mio sé”.
‘Qualsiasi percezione ci sia, passata, futura, presente, interna o
esterna, grossolana o sottile, inferiore o superiore, che sia lontana o
vicina, ogni percezione dovrebbe, per mezzo della retta saggezza, essere vista così com'è realmente:
“Questa non è mia, io non sono questo, questo non è il mio sé”.

‘Qualsiasi formazione mentale ci sia, passata, futura, presente, interna
o esterna, grossolana o sottile, inferiore o superiore, che sia lontana
o vicina, ogni formazione mentale dovrebbe, per mezzo della retta saggezza, essere vista
così com'è realmente: “Questa non è mia, io non sono questo, questo non è il
mio sé”.

\clearpage

\paliText
\markboth{\paliTitle}{\rightmark}

Taṁ kiṁ maññatha bhikkhave, viññāṇaṁ niccaṁ vā aniccaṁ vā'ti.

Aniccaṁ bhante.

Yam panāniccaṁ, dukkhaṁ vā taṁ sukhaṁ vā'ti.

Dukkhaṁ bhante.

Yam panāniccaṁ dukkhaṁ viparināma-dhammaṁ, kallaṁ nu taṁ samanupassituṁ
etaṁ mama, esoham'asmi, eso me attā'ti.

No hetaṁ bhante.

[Tasmā tiha bhikkhave] yaṁ kiñci rūpaṁ atītānāgata-paccuppannaṁ ajjhattaṁ
vā bahiddhā vā oḷārikaṁ vā sukhumaṁ vā hīnaṁ vā paṇītaṁ vā yandūre
santike vā, sabbaṁ rūpaṁ netaṁ mama, nesoham'asmi, na me so attā'ti,
evametaṁ yathābhūtaṁ sammappaññāya daṭṭhabbaṁ.

Yā kāci vedanā atītānāgata-paccuppannā ajjhattā vā bahiddhā vā oḷārikā
vā sukhumā vā hīnā vā paṇītā vā yā dūre santike vā, sabbā vedanā netaṁ
mama, nesoham'asmi, na me so attā'ti, evametaṁ yathābhūtaṁ sammappaññāya
daṭṭhabbaṁ.

Yā kāci saññā atītānāgata-paccuppannā ajjhattā vā bahiddhā vā oḷārikā vā
sukhumā vā hīnā vā paṇītā vā yā dūre santike vā, sabbā saññā netaṁ mama,
nesoham'asmi, na me so attā'ti, evametaṁ yathābhūtaṁ sammappaññāya
daṭṭhabbaṁ.

Ye keci saṅkhārā atītānāgata-paccuppannā ajjhattā vā bahiddhā vā oḷārikā
vā sukhumā vā hīnā vā paṇītā vā ye dūre santike vā, sabbe saṅkhārā netaṁ
mama, nesoham'asmi, na me so attā'ti, evametaṁ yathābhūtaṁ sammappaññāya
daṭṭhabbaṁ.

\clearpage

\englishText
\markboth{\englishTitle}{\rightmark}

‘Qualsiasi coscienza ci sia, passata, futura, presente, interna o
esterna, grossolana o sottile, inferiore o superiore, che sia lontana o vicina,
ogni coscienza dovrebbe, per mezzo della retta saggezza, essere vista così com'è realmente:
“Questa non è mia, io non sono questo, questo non è il mio sé”.

‘Vedendo in questo modo, o bhikkhu, il saggio nobile discepolo si
disincanta dalla forma, si disincanta dalla sensazione, si
disincanta dalla percezione, si disincanta dalle formazioni
mentali, si disincanta dalla coscienza. Disincantandosi,
le sue passioni svaniscono; con lo svanire della passione il
cuore è liberato; con la liberazione viene la conoscenza: “È
liberato”, ed egli sa: “La nascita è distrutta, la Santa Vita è stata
vissuta, fatto è ciò che doveva essere fatto, non c'è più ritorno a nessuno
stato dell'essere.”\thinspace ’

Così parlò il Beato. Lieti, il gruppo di cinque bhikkhu
si rallegrò di ciò che il Beato aveva detto. Inoltre, mentre questo discorso veniva
pronunciato, le menti dei cinque bhikkhu furono liberate dalle
impurità, senza più attaccamento.

Così termina il discorso sulla Caratteristica del non-sé.


\clearpage

\paliText
\markboth{\paliTitle}{\rightmark}

Yaṁ kiñci viññāṇaṁ atītānāgata-paccuppannaṁ ajjhattaṁ vā bahiddhā vā
oḷārikaṁ vā sukhumaṁ vā hīnaṁ vā paṇītaṁ vā yandūre santike vā, sabbaṁ
viññāṇaṁ netaṁ mama, nesoham'asmi, na me so attā'ti, evametaṁ yathābhūtaṁ
sammappaññāya daṭṭhabbaṁ.

[Evaṁ passaṁ bhikkhave] sutvā ariyasāvako rūpasmim pi nibbindati, vedanāya
pi nibbindati, saññāya pi nibbindati, saṅkhāresu pi nibbindati,
viññāṇasmim pi nibbindati, nibbindaṁ virajjati, virāgā vimuccati,
vimuttasmiṁ vimuttam iti ñāṇaṁ hoti, khīṇā jāti, vusitaṁ brahmacariyaṁ,
kataṁ karaṇīyaṁ, nāparaṁ itthattāyā'ti pajānātī'ti.

[Idam-avoca bhagavā.] Attamanā pañcavaggiyā bhikkhū bhagavato bhāsitaṁ
abhinanduṁ. Imasmiñca pana veyyākaraṇasmiṁ bhaññamāne pañcavaggiyānaṁ
bhikkhūnaṁ anupādāya āsavehi cittāni vimucciṁsū'ti.

Anattalakkhaṇa-suttaṁ niṭṭhitaṁ.

\chapterTocDelegatePageNumber
\chapter{The Fire Sermon}

\setTocDelegatedPageNumber
\englishText
\renewcommand{\englishTitle}{The Fire Sermon}

\begin{leader}
\soloinstr{Solo introduction}

With his skill in training the trainable, the All-transcendent Buddha,
lucid speaker, teacher of the highest knowledge,

He who expounds to the people the Dhamma and Vinaya that is fitting and
worthy, teaching with this wonderful parable about fire, meditators of
the highest skill;

He has liberated those who listen with the liberation that is utterly
complete, through true investigation, with wisdom\\ and attention.

Let us now recite this Sutta which describes the characteristics\\ of dukkha.

\end{leader}

Thus have I heard.

At one time the Blessed One was staying near Gayā at Gayā Head together
with a thousand bhikkhus. There the Blessed One addressed the bhikkhus
thus:

‘Bhikkhus, everything is burning. And what, bhikkhus, is everything
that is burning?

‘The eye, bhikkhus, is burning, forms are burning, eye consciousness is
burning, eye contact is burning, the feeling that arises from eye
contact, whether it is pleasant, painful, or neutral, that too is
burning. With what is it burning? I declare that it is burning with the
fires of passion, hatred, and delusion; it is burning with birth,
ageing, and death, with sorrow, lamentation, pain, grief, and despair.

\enlargethispage{2\baselineskip}

‘The ear is burning, sounds are burning, ear consciousness is burning,
ear contact is burning, the feeling that arises from ear contact,
whether it is pleasant, painful, or neutral, that too is burning. With
what is it burning? I declare that it is burning with the fires of
passion, hatred, and delusion; it is burning with birth, ageing, and
death, with sorrow, lamentation, pain, grief, and despair.

\chapterTocSubIndentTrue
\chapter{Āditta-pariyāya Sutta}

\paliText
\renewcommand{\paliTitle}{Āditta-pariyāya Sutta}

\begin{leader}
\soloinstr{Solo introduction}

\begin{solotwochants}
Veneyyadamanopāye  & sabbaso pāramiṁ gato\\
Amoghavacano buddho & abhiññāyānusāsako\\
Ciṇṇānurūpato cāpi & dhammena vinayaṁ pajaṁ\\
Ciṇṇāggipāricariyānaṁ & sambojjhārahayoginaṁ\\
Yamādittapariyāyaṁ & desayanto manoharaṁ\\
Te sotāro vimocesi & asekkhāya vimuttiyā\\
Tathevopaparikkhāya & viññūṇaṁ sotumicchataṁ\\
Dukkhatālakkhaṇopāyaṁ & taṁ suttantaṁ bhaṇāma se\\
\end{solotwochants}
\end{leader}

[Evaṁ me sutaṁ]

Ekaṁ samayaṁ bhagavā gayāyaṁ viharati gayāsīse saddhiṁ bhikkhu-sahassena.
Tatra kho bhagavā bhikkhū āmantesi:

Sabbaṁ bhikkhave ādittaṁ. Kiñca bhikkhave sabbaṁ ādittaṁ.

Cakkhuṁ bhikkhave ādittaṁ, rūpā ādittā, cakkhuviññāṇaṁ ādittaṁ,
cakkhusamphasso āditto, yampidaṁ cakkhusamphassapaccayā uppajjati
vedayitaṁ sukhaṁ vā dukkhaṁ vā adukkhamasukhaṁ vā tam pi ādittaṁ. Kena
ādittaṁ. Ādittaṁ rāgagginā dosagginā mohagginā, ādittaṁ jātiyā
jarāmaraṇena sokehi paridevehi dukkhehi domanassehi upāyāsehi ādittan'ti
vadāmi.

Sotaṁ ādittaṁ, saddā ādittā, sotaviññāṇaṁ ādittaṁ, sotasamphasso āditto,
yampidaṁ sotasamphassapaccayā uppajjati vedayitaṁ sukhaṁ vā dukkhaṁ vā
adukkhamasukhaṁ vā tam pi ādittaṁ. Kena ādittaṁ. Ādittaṁ rāgagginā
dosagginā mohagginā, ādittaṁ jātiyā jarāmaraṇena sokehi paridevehi
dukkhehi domanassehi upāyāsehi ādittan'ti vadāmi.

\clearpage

\englishText
\markboth{\englishTitle}{\rightmark}

‘The nose is burning, odours are burning, nose consciousness is burning,
nose contact is burning, the feeling that arises from nose contact,
whether it is pleasant, painful, or neutral, that too is burning. With
what is it burning? I declare that it is burning with the fires of
passion, hatred, and delusion; it is burning with birth, ageing, and
death, with sorrow, lamentation, pain, grief, and despair.

‘The tongue is burning, tastes are burning, tongue consciousness is
burning, tongue contact is burning, the feeling that arises from tongue
contact, whether it is pleasant, painful, or neutral, that too is
burning. With what is it burning? I declare that it is burning with the
fires of passion, hatred, and delusion; it is burning with birth,
ageing, and death, with sorrow, lamentation, pain, grief, and despair.

‘The body is burning, tangible objects are burning, body consciousness
is burning, body contact is burning, the feeling that arises from body
contact, whether it is pleasant, painful, or neutral, that too is
burning. With what is it burning? I declare that it is burning with the
fires of passion, hatred, and delusion; it is burning with birth,
ageing, and death, with sorrow, lamentation, pain, grief, and despair.

‘The mind is burning, mental states are burning, mind consciousness is
burning, mind contact is burning, the feeling that arises through mind
contact, whether it is pleasant, painful, or neutral, that too is
burning. With what is it burning? I declare that it is burning with the
fires of passion, hatred, and delusion; it is burning with birth,
ageing, and death, with sorrow, lamentation, pain, grief, and despair.

‘Seeing thus, bhikkhus, the wise noble disciple becomes disenchanted
with the eye, disenchanted with forms, disenchanted with eye
consciousness, disenchanted with eye contact, and the feeling that
arises from eye contact --- whether it is pleasant, painful, or
neutral --- that too they become disenchanted with.

\clearpage

\paliText
\markboth{\paliTitle}{\rightmark}

Ghānaṁ ādittaṁ, gandhā ādittā, ghānaviññāṇaṁ ādittaṁ, ghānasamphasso
āditto, yampidaṁ ghānasamphassapaccayā uppajjati vedayitaṁ sukhaṁ vā
dukkhaṁ vā adukkhamasukhaṁ vā tam pi ādittaṁ. Kena ādittaṁ. Ādittaṁ
rāgagginā dosagginā mohagginā, ādittaṁ jātiyā jarāmaraṇena sokehi
paridevehi dukkhehi domanassehi upāyāsehi ādittan'ti vadāmi.

Jivhā ādittā, rasā ādittā, jivhāviññāṇam ādittaṁ, jivhāsamphasso āditto,
yampidaṁ jivhāsamphassapaccayā uppajjati vedayitaṁ sukhaṁ vā dukkhaṁ vā
adukkhamasukhaṁ vā tam pi ādittaṁ. Kena ādittaṁ. Ādittaṁ rāgagginā
dosagginā mohagginā, ādittaṁ jātiyā jarāmaraṇena sokehi paridevehi
dukkhehi domanassehi upāyāsehi ādittan'ti vadāmi.

Kāyo āditto, phoṭṭhabbā ādittā, kāyaviññāṇaṁ ādittaṁ, kāyasamphasso
āditto, yampidaṁ kāyasamphassapaccayā uppajjati vedayitaṁ sukhaṁ vā
dukkhaṁ vā adukkhamasukhaṁ vā tam pi ādittaṁ. Kena ādittaṁ. Ādittaṁ
rāgagginā dosagginā mohagginā, ādittaṁ jātiyā jarāmaraṇena sokehi
paridevehi dukkhehi domanassehi upāyāsehi ādittan'ti vadāmi.

Mano āditto, dhammā ādittā, manoviññāṇaṁ ādittaṁ, manosamphasso āditto,
yampidaṁ manosamphassapaccayā uppajjati vedayitaṁ sukhaṁ vā dukkhaṁ vā
adukkhamasukhaṁ vā tam pi ādittaṁ. Kena ādittaṁ. Ādittaṁ rāgagginā
dosagginā mohagginā, ādittaṁ jātiyā jarāmaraṇena sokehi paridevehi
dukkhehi domanassehi upāyāsehi ādittan'ti vadāmi.

\enlargethispage{2\baselineskip}

[Evaṁ passaṁ bhikkhave] sutvā ariyasāvako cakkhusmiṁ pi nibbindati,
rūpesu pi nibbindati, cakkhuviññāṇe pi nibbindati, cakkhusamphassepi
nibbindati, yampidaṁ cakkhusamphassapaccayā uppajjati vedayitaṁ sukhaṁ
vā dukkhaṁ vā adukkhamasukhaṁ vā tasmiṁ pi nibbindati.

\clearpage

\englishText
\markboth{\englishTitle}{\rightmark}

‘They become disenchanted with the ear, disenchanted with sounds,
disenchanted with ear consciousness, disenchanted with ear contact, and
the feeling that arises from ear contact --- whether it is pleasant,
painful, or neutral --- that too they become disenchanted with.

‘They become disenchanted with the nose, disenchanted with odours,
disenchanted with nose consciousness, disenchanted with nose contact,
and the feeling that arises from nose contact --- whether it is pleasant,
painful, or neutral --- that too they become disenchanted with.

‘They become disenchanted with the tongue, disenchanted with tastes,
disenchanted with tongue consciousness, disenchanted with tongue
contact, and the feeling that arises from tongue contact --- whether it is
pleasant, painful, or neutral --- that too they become disenchanted with.

‘They become disenchanted with the body, disenchanted with tangible
objects, disenchanted with body consciousness, disenchanted with body
contact, and the feeling that arises from body contact --- whether it is
pleasant, painful, or neutral --- that too they become disenchanted with.

‘They become disenchanted with the mind, disenchanted with mental
states, disenchanted with mind consciousness, disenchanted with mind
contact, and the feeling that arises from mind contact --- whether it is
pleasant, painful, or neutral --- that too they become disenchanted with.

‘Becoming disenchanted, their passions fade away; with the fading of
passion the heart is liberated; with liberation there comes the
knowledge: “It is liberated,” and they know: “Destroyed is birth, the
Holy Life has been lived out, done is what had to be done, there is no
more coming into any state of being.”\thinspace ’

\enlargethispage{\baselineskip}

Thus spoke the Blessed One; delighted, the bhikkhus rejoiced in what the
Blessed One had said. Moreover, while this discourse was being uttered, the
minds of those thousand bhikkhus were freed from the defilements,
without any further attachment.

Thus ends The Fire Sermon.

\clearpage

\paliText
\markboth{\paliTitle}{\rightmark}

Sotasmiṁ pi nibbindati, saddesu pi nibbindati, sotaviññāṇe pi
nibbindati, sotasamphassepi nibbindati, yampidaṁ sotasamphassapaccayā
uppajjati vedayitaṁ sukhaṁ vā dukkhaṁ vā adukkhamasukhaṁ vā tasmiṁ pi
nibbindati.

Ghānasmiṁ pi nibbindati, gandhesu pi nibbindati, ghānaviññāṇe pi
nibbindati, ghānasamphassepi nibbindati, yampidaṁ ghānasamphassapaccayā
uppajjati vedayitaṁ sukhaṁ vā dukkhaṁ vā adukkhamasukhaṁ vā tasmiṁ pi
nibbindati.

Jivhāya pi nibbindati, rasesu pi nibbindati, jivhāviññāṇe pi nibbindati,
jivhāsamphassepi nibbindati, yampidaṁ jivhāsamphassapaccayā uppajjati
vedayitaṁ sukhaṁ vā dukkhaṁ vā adukkhamasukhaṁ vā tasmiṁ pi nibbindati.

Kāyasmiṁ pi nibbindati, phoṭṭhabbesu pi nibbindati, kāyaviññāṇe pi
nibbindati, kāyasamphassepi nibbindati, yampidaṁ kāyasamphassapaccayā
uppajjati vedayitaṁ sukhaṁ vā dukkhaṁ vā adukkhamasukhaṁ vā tasmiṁ pi
nibbindati.

Manasmiṁ pi nibbindati, dhammesu pi nibbindati, manoviññāṇe pi
nibbindati, manosamphassepi nibbindati, yampidaṁ manosamphassapaccayā
uppajjati vedayitaṁ sukhaṁ vā dukkhaṁ vā adukkhamasukhaṁ vā tasmiṁ pi
nibbindati.

Nibbindaṁ virajjati, virāgā vimuccati, vimuttasmiṁ, vimuttam iti ñāṇaṁ
hoti, khīṇā jāti, vusitaṁ brahmacariyaṁ, kataṁ karaṇīyaṁ, nāparaṁ
itthattāyā'ti pajānātī'ti.

\enlargethispage{\baselineskip}

[Idam-avoca bhagavā.] Attamanā te bhikkhū bhagavato bhāsitaṁ abhinanduṁ.
Imasmiñca pana veyyākaraṇasmiṁ bhaññamāne tassa bhikkhu-sahassassa
anupādāya āsavehi cittāni vimucciṁsū'ti.

Ādittapariyāya-suttaṁ niṭṭhitaṁ.

\resumeNormalText

% End of suttas.tex
